% RESEARCH DRAFT — do not include in main.tex until Lemmas 2--4 and
% Proposition~\ref{prop:clipped-calibration} have passed independent audit.
% Relies on the specific uniform bounded-observable reduction of Lemma 2.
\subsection*{Localising the clipping tail to exceptional mass}
\label{subsec:exceptional-mass}
Retain the definitions of Proposition~\ref{prop:clipped-calibration}:
\[
 W_X=(X^{1-\beta},X]\cap\mathbb N,\quad
 \beta=1/200,\quad H_X=X(\log X)^{\varepsilon_0/2},
 \quad B_X=\sum_{a\in W_X}a\ll X^2.
\]
For each fixed $X$, let $\mu_X$ denote the logarithmic limiting
probability distribution of $\Phi_X(w^\flat)$ as $w$ varies over the
positive integers with weights $1/w$. This distribution exists because
$\Phi_X$ has finite range in $[0,B_X]$ and each indicator
$\mathbf1\{\Phi_X\in A\}$ satisfies the first-passage invariance
hypothesis of Lemma~\ref{lem:reduction}. Let $\nu_X$ be the law of
$\Phi_X(\mathbf N)$ for the logarithmically sampled odd starting value
$\mathbf N\sim\mathbf{Log}_{u,v}$ in the proof of
Proposition~\ref{prop:clipped-calibration}.

\begin{lemma}[Distribution-level reduction]
\label{lem:distributional-reduction}
Conditional on the uniform bounded-observable estimate in
Lemma~\ref{lem:reduction},
\[
 d_{\mathrm{TV}}(\mu_X,\nu_X)
 \ll(\log X)^{-\varepsilon_0}.
\]
\end{lemma}
\begin{proof}
For an arbitrary subset $A\subseteq[0,B_X]$, apply
Lemma~\ref{lem:reduction} to
$\varphi_A(m)=\mathbf1\{\Phi_X(m)\in A\}$. The invariance required
there is exactly the first-passage invariance proved in Step~2 of
Theorem~\ref{thm:lower}, and the error constant does not depend on $A$.
Taking the supremum of the probability differences over $A$ proves
the total-variation claim.
\end{proof}

\begin{proposition}[Exceptional-mass localisation; conditional]
\label{prop:exceptional-mass}
Assume Lemmas~\ref{lem:reduction}--\ref{lem:model}, the coupling in
Step~3 of Theorem~\ref{thm:lower}, and the full ideal-window second-moment
bound $\mathbb E(Z_X^+)^2\ll X^2$ of
Lemma~\ref{lem:full-ideal-second-moment} in the separate working
addendum on intermediate tails.
There exists a coupling of $Y_X\sim\mu_X$ with the ideal-window
variable $Z_X^+$ and an event $E_X$ such that
\begin{equation}\label{eq:exception-probability}
 \mathbb P(E_X)\ll\delta_X,\qquad
 \delta_X:=(\log X)^{-\varepsilon_0}+X^{-\eta},
\end{equation}
and, on $E_X^c$,
\[
 Y_X\le(1+O(X^{-1/2}))Z_X^+.
\]
Moreover, for every $H\ge X$,
\begin{equation}\label{eq:exceptional-tail-bound}
 -C H\delta_X
 \le \mathbb E(Y_X-H)_+-\mathbb E[Y_X\mathbf1_{E_X}]
 \le C\frac{X^2}{H},
\end{equation}
where the upper bound has an absolute implied constant depending only on
fixed parameters. In particular, for $H=H_X$,
\begin{equation}\label{eq:exceptional-tail-asymptotic}
 \boxed{T_X=\mathbb E[Y_X\mathbf1_{E_X}]+o(X).}
\end{equation}
Consequently $T_X=o(X)$ if and only if
$\mathbb E[Y_X\mathbf1_{E_X}]=o(X)$ for this coupling.
\end{proposition}
\begin{proof}
The preceding lemma and maximal coupling give a random variable
$Y_X\sim\mu_X$ on an extension of the probability space of
$\Phi_X(\mathbf N)$ such that
$\mathbb P(Y_X\ne\Phi_X(\mathbf N))\ll(\log X)^{-\varepsilon_0}$.
Retain the joint distribution of $\mathbf N,Z_X^+$ and the successful
valuation coupling event $G_X$; declare
$E_X=\{Y_X\ne\Phi_X(\mathbf N)\}\cup G_X^c$.
The desired pointwise comparison holds off $E_X$, and
\eqref{eq:exception-probability} follows.

On $E_X^c$, the elementary inequality
$(y-H)_+\le y^2/H$ for $y\ge0$, together with the model second moment,
gives
\[
 \mathbb E[(Y_X-H)_+\mathbf1_{E_X^c}]
 \ll X^2/H.
\]
On $E_X$, $0\le Y_X-(Y_X-H)_+=\min(Y_X,H)\le H$. Thus
\[
 \mathbb E[Y_X\mathbf1_{E_X}]-H\mathbb P(E_X)
 \le\mathbb E[(Y_X-H)_+\mathbf1_{E_X}]
 \le\mathbb E[Y_X\mathbf1_{E_X}].
\]
Combining the two bounds proves \eqref{eq:exceptional-tail-bound}.
For $H=H_X$, both $X^2/H_X$ and $H_X\delta_X$ are $o(X)$.
By the definition of $T_X$ in Corollary~\ref{cor:two-tails},
$T_X=\mathbb E(Y_X-H_X)_+$, proving the conclusion.
\end{proof}

\begin{corollary}[Intrinsic high-mass criterion; conditional]
\label{cor:intrinsic-exceptional-mass}
Keep the preceding assumptions and put
\[
 K_X=\frac{X(\log X)^{\varepsilon_0}}{\log\log X}.
\]
For all sufficiently large $X$,
\[
 \mu_X\{Y_X>K_X\}\ll\frac{X^2}{K_X^2}+\delta_X,
 \qquad K_X\mu_X\{Y_X>K_X\}=o(X),
\]
and
\begin{equation}\label{eq:intrinsic-exceptional-mass}
 \boxed{T_X=\mathbb E_{\mu_X}
 \bigl[Y_X\mathbf1_{\{Y_X>K_X\}}\bigr]+o(X).}
\end{equation}
Hence $T_X=o(X)$ is equivalent to the \emph{intrinsic} condition
\[
 \mathbb E_{\log}\bigl[
 \Phi_X(w^\flat)\mathbf1_{\{\Phi_X(w^\flat)>K_X\}}
 \bigr]=o(X).
\]
\end{corollary}
\begin{proof}
On the good coupling event, Chebyshev's inequality for $Z_X^+$
gives the first estimate. Multiplication by $K_X$ yields
$O(X^2/K_X+K_X\delta_X)=o(X)$.
The intermediate-tail proposition
\ref{prop:intermediate-tail} gives
$\mathbb E(Y_X-H_X)_+
 =\mathbb E(Y_X-K_X)_++o(X)$, while
\[
 \mathbb E[Y_X\mathbf1_{Y_X>K_X}]
 =\mathbb E(Y_X-K_X)_++K_X\mu_X\{Y_X>K_X\}.
\]
The last term is $o(X)$, proving the claim.
\end{proof}

\begin{corollary}[An unweighted return-count criterion]
\label{cor:return-count}
For odd $m$, let $V_X(m)$ be the number of target points in $W_X$
visited in the Syracuse blocks through first passage below $X_1$,
with $V_X(m)=0$ if that passage does not occur. If for some
$\theta>0$ the actual logarithmic starting-value law satisfies
\begin{equation}\label{eq:return-count-moment}
 \sup_{X\ge X_0}\mathbb E_{\log}
     V_X(w^\flat)^{1+\theta}<\infty,
\end{equation}
then $T_X=o(X)$. More precisely,
\[
 T_X\ll X(\log X)^{-\varepsilon_0\theta/2}.
\]
\end{corollary}
\begin{proof}
The sum defining $\Phi_X$ contains exactly $V_X$ target entries,
each at most $X$, so $0\le\Phi_X\le XV_X$. If
$\Phi_X>H_X$, then $V_X>H_X/X$. Thus
\[
 T_X=\mathbb E_{\log}(\Phi_X-H_X)_+
 \le X\mathbb E_{\log}\bigl[V_X\mathbf1_{V_X>H_X/X}\bigr]
 \le X(H_X/X)^{-\theta}\mathbb E_{\log}V_X^{1+\theta},
\]
proving the assertion. The missing estimate is
\eqref{eq:return-count-moment} for \emph{actual} Syracuse orbits;
Lemma~\ref{lem:model} only controls the corresponding ideal model.
\end{proof}

\begin{remark}[A sharp limitation of the available estimates]
\label{rem:exception-mass-obstruction}
The probability bound in \eqref{eq:exception-probability} by itself
cannot give $T_X=o(X)$. For illustration, consider an \emph{abstract}
random variable (not an actual Collatz orbit) equal to
$\gamma_0X$ except on an event of probability
$(\log X)^{-\varepsilon_0}$, where it equals $X^2/4$.
It is within that total-variation distance of a model concentrated at
$\gamma_0 X$, and its clipped mean at $H_X$ equals
$\gamma_0X+o(X)$, yet its unclipped tail is asymptotic to
$X^2/[4(\log X)^{\varepsilon_0}]$, which is not $o(X)$.
The remaining task is therefore an $L^1$ estimate on exceptional
\emph{actual-orbit mass}, not simply a smaller upper bound for the
frequency of atypical valuation strings.
\end{remark}
